Figure~\ref{fig:result_widths} shows the final measured widths of 
  the away-side correlation (standard-deviation of $Y^{\mathrm{corr}}(\Dphi)$ in Eq.~\eqref{eq:fit_func})
  in \PbPb\ events, as a function of centrality.
The value measured in \pp\ collisions is also shown.
The measurements are shown for All pairs, Same-sign pairs 
  and Opposite-sign pairs.
The widths of the same- and opposite-sign pairs are similar 
  with each other for both the \PbPb\ and \pp\ measurements.
The measured \PbPb\ values are very similar to the values measured in 
  \pp\ collisions.
No clear increase in the widths as a function of centrality 
  is observed.
Instead a slight decrease in the 0--10\% most central events is observed.
Figure~\ref{fig:result_widths_fits} shows linear fits to the
  \PbPb\ widths as a function of centrality.
The slopes are negative for all three charge-combinaitons
  but consistent with zero within 3$\sigma$ statistical uncertainties.
Figure~\ref{fig:result_widths_diff} shows the squared difference
  between the \pp\ and \PbPb\ measurements, which further quantifies
  the difference between the \pp\ and \PbPb\ measurements.
In several centrality intervals the \pp\ widths are larger
  than the \PbPb\ values, though always within 1--2~$\sigma$
  statistical/systematic uncertainties.
This is very surprizing as one would expect that interactions
  with the QGP would broaden the \Dphi\ correlations~\cite{Nahrgang:2013saa}.
This unexpected behavior is investigated next.

\begin{figure}
\centering
\includegraphics[width=1.0\textwidth]%                                                                                           
{figures/AllFigs/can_systPbPbPP_WidthTau_Tight_pT_GT4_EffCor}
\caption{
The measured widths of the away-side correlation.
The plot shows the $\sigma$ (standard-deviation of $Y^{\mathrm{corr}}(\Dphi)$ in Eq.\eqref{eq:fit_func})
  as well as the $\Gamma$ (the half-width at half-maximum parameter in Eq.\eqref{eq:fit_func}))
  for dimuon pairs, as a function of centrality.
Also shown for comparison, are the corresponding parameters in \pp\ collisions (solid markers).
  The horizontal lines indicates the nominal \pp\
  values (i.e. the value central value of the \pp\ measurements,
  but without the statistical or systematic uncertainties) drawn across the full centrality range.
The left, middle and right panels correspond to
  Same-sign pairs, Opposite-sign pairs and All-pairs.
The vertical lines and shaded bars indicate statistical and
  systematic uncertainties, respectively.
\label{fig:result_widths}
}
\end{figure}



\begin{figure}
\centering
\includegraphics[width=1.0\textwidth]%
{figures/AllFigs/can_fits_0_Tight_pT_GT4_EffCor}
\caption{
\label{fig:result_widths_fits}
Linear fits to the measured widths of the away-side correlation.
The left, middle and right panels correspond to
  Same-sign pairs, Opposite-sign pairs and All-pairs.
The top row shows $\chi^2$-fits  made taking into consideration the statistical uncertainties only.
The middle (bottom) row shows $\chi^2$-fits  made taking into consideration the systematic  uncertainties only
  while treating the uncertainties fully uncorrelated (correlated) bin-to-bin.
The fit values and their uncertainties are shown on the individual panels.
}
\end{figure}

\begin{figure}
\centering
\includegraphics[width=1.0\textwidth]%
{figures/AllFigs/can_Rmsdiff_0_Tight_pT_GT4_EffCor}
\caption{
The difference between the squared-widths of the \pp\ and \PbPb\ measurements.
The left, middle and right panels correspond to
  Same-sign pairs, Opposite-sign pairs and All-pairs.
The vertical lines and shaded bars indicate statistical and
  systematic uncertainties, respectively.
\label{fig:result_widths_diff}
}
\end{figure}


Figure~\ref{fig:ptbar_tight} shows the mean-\ptbar\ values, 
  where $\ptbar=(\pt^{\mu_a}+\pt^{\mu_b})/2$,
  for the muon-pairs contributing to the pair distributions.
The mean values are calculated for the difference of the $\ptbar$
  distributions on the away-side ($|\Dphi|<\pi2$)
  and near-side ($|\Dphi|>\pi/2$).
The \ptbar\ distributions also
  have the restriction that $\ptbar<20\GeV$ and that
  the \pt\ of the individual muons are greater than 4~\GeV.
It is observed that for peripheral events (60--80\% centrality),
  the \pp\ and \PbPb\ $\langle\ptbar\rangle$ values are in agreement.
For more central collisions the $\langle\ptbar\rangle$ values
  for the \PbPb\ data are higher,especially for the opposite-sign case.
Figure~\ref{fig:ptbar_medium} shows similar results for
  the Medium muon selection, which gives values consistent 
  with the Tight selection.
This hardening of the \ptbar\ spectras, from peripheral
  to central collisions, can arises from multiple effects 
  such as:
\begin{itemize}
  \item Stronger suppression for parent quarks of lower-\pt\ muons
      in more central collisions.
  \item Different relative contributions from Drell-Yan processes 
    (for the opposite-sign pairs) as a function of centrality.
%  \item Different rates of suppression of $b$ vs $c$ quarks as 
%	  function of centrality.
%	This would also indicate why the $\langle\ptbar\rangle$
%	  shows a smaller variation in the same-sign case,
%	  which is almost exclusively from $\bbar$ pairs.
\end{itemize}
The harder spectra in more central collisions, leads
  to a narrowing of the \Dphi-widths of the correlations,
  as the widths decrease with increasing \pt.
Next we do a semi-quantitative study of whether the increase 
  in the $\langle\ptbar\rangle$ (Figure~\ref{fig:ptbar_tight}) 
  from peripheral to central collisions, is sufficient to produce
  the decrease in the widths for more central collisions.
Figure~\ref{fig:ptbar_tight} also shows the 
  \pp-$\langle\ptbar\rangle$ values for different cuts 
  on the minimum \ptbar.
The increase in the $\langle\ptbar\rangle$ 
  in the 0--10\% most central \PbPb\ collisions, 
  in the Same-sign case can be reproduced by increasing 
  the lower-cut on the \ptbar\ in the \pp\ data to 5~\GeV.
While the increase in the opposite-sign case can be reproduced
  by increasing the lower-cut on the \ptbar\ in the \pp\ data
  to 6.5~\GeV.
With these additional cuts on \ptbar, the widths of the correlations
  are re-measured for the \pp\ data, and are compared to the nominal-\PbPb\
  measurements in Figure~\ref{fig:widths_comparepp_multiplepT}.
It is seen that increasing the \pp-$\langle\ptbar\rangle$
  is sufficient to reproduce the decrease in the widths
  observed in the \PbPb\ data.
This study then indicates that the effect of the suppression 
  of HF-particles, which increases the $\langle\ptbar\rangle$
  leads to an decrease in the measured widths from 
  peripheral$\rightarrow$central collisions.
The effects of angular decorrelation of the HF-pair, if any,
  are supersceded by the effects of the decrease in the
  widths that arise from the hardening of the spectra.
  

\begin{figure}[h] 
\centering
\includegraphics[width=1.0\textwidth]%
{figures/AllFigs/can_pTBar_option1_Tight_pT_GT4_EffCor}
\caption{
The$\langle\ptbar\rangle$ values for the muon-pairs
  contributing to the pair distributions.
The mean values are calculated for the difference of the $\ptbar$
  distributions on the away-side ($|\Dphi|<\pi2$) 
  and near-side ($|\Dphi|>\pi/2$), with the restriction
  that $4<\pt^{\mu}/\GeV<20$.
The \PbPb\ values are shown for different centrality intervals ($x$-axis).
The horizontal line indicates the value for the \pp\ data.
The left, center and right panels correspond to same-charge,
  opposite-charge and combined-charge pairs.
\label{fig:ptbar_tight}
}
\end{figure}



\begin{figure}[h] 
\centering
\includegraphics[width=1.0\textwidth]%
{figures/AllFigs/can_systPbPbPPMultiplepT_0_Tight_pT_GT4_EffCor}
\caption{
Comparison of the widths in \PbPb\ collisions (from Figure~\ref{fig:result_widths})
  to those measured in \pp\ collisions with different
  cuts on the \ptbar.
\label{fig:widths_comparepp_multiplepT}
}
\end{figure}






\iftoggle{IncludeYields}
{
  Figure~\ref{fig:result_yields}, shows the measured yields of 
    the away-side correlation (integral of $Y^{\mathrm{corr}}(\Dphi)$ in Eq.\eqref{eq:fit_func})
    in \PbPb\ events, as a function of centrality.
  Figure~\ref{fig:result_yields_taa}, shows similar plots, but with for the
    per-event yields scaled by the \TAA (Eq.~\eqref{eq:TaaYield}).
  The cross-section measured in \pp\ collisions is also shown in Figure~\ref{fig:result_yields_taa}.
  As expected, the yields (Figure~\ref{fig:result_yields}) increase monotonically
    from peripheral to central events.
  However the \TAA\ scaled yields show a monotonic decrease. 
  This indicates that there is a significant suppression in the number of dimuon pairs
    compared to what one would expect if the \PbPb\ collisions were
    a superimposition of independent \pp\ collisions; 
    in which case the scaled yields would be a constant and consistent 
    with the cross-section measured in \pp\ collisions.
  The relative decrease in the \TAA\ scaled-yields are similar for the same-sign pairs 
    and the opposite-sign pairs.

  \begin{figure}
  \centering
  \includegraphics[width=1.0\textwidth]%                                                                                           
  {figures/AllFigs/can_syst_1_Tight_pT_GT4_EffCor}
  \caption{
  The measured correlated yields on the away-side 
    (integral of $Y^{\mathrm{corr}}(\Dphi)$ in Eq.\eqref{eq:fit_func}) for dimuon pairs, 
    as a function of centrality.
  The left, middle and right panels correspond to
    All-pairs, Same-sign pairs and Opposite-sign pairs.
  The vertical lines and shaded bars indicate statistical and 
    systematic uncertainties, respectively.
  The overall luminosity uncertainty of $\pm1.5$\% is not included
    in the shaded bands, instead stated on the plot.
  \label{fig:result_yields}
  }
  \end{figure}

  \begin{figure}
  \centering
  \includegraphics[width=1.0\textwidth]%                                                                                           
  {figures/AllFigs/can_systPbPbPP_2_Tight_pT_GT4_EffCor}
  \caption{
  The \TAA\ scaled yields on the away-side for dimuon pairs, 
    as a function of centrality.
  Also shows for comparison, is the \pp\ cross-section (blue marker).
  The horizontal blue-line indicates the nominal \pp\
    value drawn across the full centrality range.
  The left, middle and right panels correspond to
    All-pairs, Same-sign pairs and Opposite-sign pairs.
  The vertical lines and shaded bars indicate statistical and 
    systematic uncertainties, respectively.
  The overall luminosity uncertainties of $\pm1.5$\% for \PbPb\
    and $\pm1.6$\% for \pp\ are not included
    in the shaded bands, instead stated on the plot.
  \label{fig:result_yields_taa}
  }
  \end{figure}
}
{}

\clearpage
\section{Summary}
This note reports a measurement of correlated back-to-back 
  muon-pair production in \PbPb\ collisions at \snn=5.02~\TeV. 
  and in \pp\ collisions at \sqs=5.02~\TeV.
The \PbPb\ measurements are performed using data from the 2015 and 2018
  heavy-ion runs corresponding to an integrated luminosity of
  1.93~\inb.
While the \pp\ measurements are performed using data from the 2017
  reference run corresponding to an integrated luminosity of 0.26~\ifb.
The measurements are performed over the integrated \pt>4~\GeV\ range.
Over this \pt\ range, nearly all the dimuon pairs arise from decay
  of HF hadrons.
Backgrounds from the underlying event (including flow-correlated pairs) 
  are removed.
Backgrounds from muons arising from inflight decays of pions and Kaons,
  as well as muons arising from secondary interactions in the calorimeter 
  were estimated using the momentum-imbalance as a discriminant, 
  and were found to be below a few percent.
\iftoggle{IncludeYields}
{
  The yields of the correlated muon-pairs, as well as the width of the
    correlation in \Dphi\ are measured.
  The per-event yields, scaled by the nuclear thickness function, \TAA, 
    are observed to decrease from peripheral to central collisions,
    indicating a suppression in the yield of such pairs in central collisions.
  On the other hand, the widths of the correlations are observed to be 
    quite independent of centrality.
  Theoretical models~\cite{Nahrgang:2013saa} predict varying amounts of 
    suppression and broadening of back-to-back HF pairs depending 
    on the energy-loss mechanism.
  Thus these measurements can shed light on the energy-loss mechanism of 
    HF-quarks in the QGP produced in heavy-ion collisions.
}
{
  The width of the correlated muon-pairs in \Dphi\ are measured,
    and are observed to be independent of centrality, within uncertainties.
  The widths are also observed to be consistent between the \pp\ and the \PbPb\ systems.
  Theoretical models~\cite{Nahrgang:2013saa} predict varying amounts of 
    centrality-dependent broadening of such back-to-back HF pairs in \PbPb\ collisions
    depending on the heavy-quark energy-loss mechanism.
  Thus these measurements can shed light on the energy-loss mechanism of 
    HF-quarks in the QGP produced in heavy-ion collisions.
}
